\documentclass[11pt]{article}
\usepackage[margin=1in]{geometry}
\usepackage{amsmath,amssymb,amsthm}
\usepackage{booktabs}
\usepackage{enumitem}
\usepackage[hidelinks]{hyperref}
\newtheorem{definition}{Definition}
\newtheorem{remark}{Remark}
\newcommand{\Cmap}{C}
\newcommand{\Dmap}{D}

\title{Understanding as a Lossless Round-Trip Fixed Point:\\
A Descriptive Formalization of General Thinking with an Injectable Protocol}
\author{[Author]}
\date{}

\begin{document}
\maketitle

\begin{abstract}
We present a descriptive formalization of ``general'' thinking---reasoning that
moves beyond a single domain---organized as a compression--decomposition round
trip. We define understanding as a \emph{lossless round-trip fixed point}: a
compression $\Cmap$ from phenomena to models and a decompression $\Dmap$ from
models to predictions such that $\Dmap\circ\Cmap$ acts as the identity on the
domain of interest. Losslessness is taken to mean uniqueness plus reproducibility
(equivalently, controlled predictive quality). From this single condition we
derive four facets (predictive power, perturbation, fitness, leading-term
anchoring) and two derivations (basis choice via multiresolution analysis, and a
two-stage domain-selection-then-refinement procedure whose order cannot be
reversed). We test the framework against seven thinking-style hypotheses drawn
from the source corpus and obtain 7/7 internal consistency (explicitly an
internal consistency check on the source corpus, not external validation). We
finally distill five injectable prompt templates (T1--T5), each with a trigger
condition and an action, and position them as an agent-prompt injection protocol.
We claim injectability and operability only; effectiveness is left to independent
experiments.
\end{abstract}
\section{Introduction}
\label{sec:intro}

Why can ``how to think generally'' be formalized at all? General thinking---
moving from a concrete solution to a cost model to a unified form, checking that
the original solution is a special case, and landing every qualitative claim on
a testable number---looks like a craft. We argue that it has a stable internal
structure that admits a \emph{descriptive} formalization: a small set of
principles from which the observed practices follow as consequences. Our
contribution is descriptive; it is not a claim about the mechanisms of human or
artificial cognition.

We make three contributions. \emph{First}, a core formalization: understanding
as a lossless round-trip fixed point of a compression--decomposition pair
($\Dmap\circ\Cmap=\mathrm{id}$), with losslessness defined as uniqueness plus
reproducibility. \emph{Second}, a consistency test: we derive seven
thinking-style hypotheses from the framework and obtain 7/7 internal consistency
on the source corpus (an internal consistency check, not external validation).
\emph{Third}, an actionable protocol: five injectable prompt templates
(T1--T5), each specifying a trigger condition and an action, positioned as an
agent-prompt injection protocol.

\paragraph{Scope and honesty.}
This paper is a descriptive formalization plus an actionable protocol. We do
\emph{not} claim to have discovered the mechanism of human or AI thinking, nor
do we offer cognitive-science conclusions. The consistency test is
self-referential: the framework was constructed from and tested on the same
corpus. Template effectiveness is pending independent experiments; we claim
injectability and operability only. All quantitative statements (7/7, five
templates, four facets) are traceable to the source corpus (see the appendix and
the accompanying number-trace file).
\section{Core Formalization: Understanding as a Lossless Round-Trip Fixed Point}
\label{sec:core}

\begin{definition}[Lossless round-trip fixed point]
\label{def:fixedpoint}
A model $m$ of a phenomenon space $\mathcal{X}$ constitutes \emph{understanding}
if there exist a compression $\Cmap:\mathcal{X}\to\mathcal{M}$ (phenomena to
model) and a decompression $\Dmap:\mathcal{M}\to\mathcal{X}$ (model to
phenomena/predictions) such that $\Dmap\circ\Cmap = \mathrm{id}$ on the domain of
interest. A compression state satisfying this condition is called a
\emph{lossless round-trip fixed point}.
\end{definition}

\paragraph{Losslessness.}
We take ``lossless'' to mean \emph{unique} plus \emph{reproducible}: the
compressed state decompresses back to the same phenomenon (up to controlled
tolerance) on every attempt. In the source corpus this is phrased equivalently
as \emph{predictive quality}: a compressed state that uniquely and reproducibly
expands back to the phenomena carries the predictive power that losslessness
demands. Losslessness is therefore a property of the round trip, not of the
compression ratio alone.

\paragraph{Self-reference is finite.}
Finding the fixed point is a self-referential task (the model must account for
its own compression), but the fixed-point equation $x' = x$ is typically
\emph{finite}: it is solved by iteration converging to a solution (as in solving
$Ce^{t}$-type recursions), not an infinite regress. Pathological self-reference
(G\"odel-type incompleteness) applies only to degenerate self-referential
systems; ordinary modeling fixed points are well-posed. This finiteness is what
makes the definition operational rather than circular.

\paragraph{Reading the definition.}
The definition is deliberately weak: it names a condition, not a mechanism. Its
role is organizational---every practice we describe below (prediction,
perturbation, fitness, anchoring, basis choice, two-stage compression) is a
consequence of demanding the round trip close. Where the round trip cannot be
closed, the framework says the model is not yet understood, and points to which
facet failed (Section~\ref{sec:facets}).
\section{Four Facets of the Round Trip}
\label{sec:facets}

The single condition $\Dmap\circ\Cmap=\mathrm{id}$ has four faces. Each facet is
a way the round trip can be examined, stressed, or made to fail.

\paragraph{1. Predictive power (temporal direction).}
Predictive power is the round trip executed in the temporal direction: expanding
from the compressed state forward into future phenomena. The error is
controlled and nonzero---a perfect identity on the past is neither expected nor
required; what matters is that the expansion's error stays within a known bound.
In the source corpus this is the face used to grade models: a model earns its
keep by expanding into testable, falsifiable predictions.

\paragraph{2. Perturbation (how the round trip is achieved).}
Perturbation theory is the \emph{implementation} of the round trip: represent
the decompression as a leading term plus successive corrections, and the round
trip closes within a convergence radius ($\Dmap\circ\Cmap\approx\mathrm{id}$).
Necessity enters as a rule: higher-order terms are added only when the leading
terms fail to close the round trip. This turns ``how much detail to keep'' into
a measurable question (which correction term first exceeds tolerance), not a
taste judgment.

\paragraph{3. Fitness (the landscape of thinking).}
There is no binary ``wrong''; there is a continuous \emph{fitness} measure on
candidate models. Progress is gradient descent on this landscape: incremental,
local moves are robust to noise, which is why stepwise refinement outperforms
occasional wholesale replacement. Fitness is the facet that makes the framework
compatible with empirical work: it is the quantity that experiments can measure.

\paragraph{4. Leading-term anchoring (robustness to randomness).}
Robustness to randomness comes from anchoring on the leading term: the leading
term occupies the widest valley in the fitness landscape, so it is the safest
thing to grasp first; perturbations are then handled term by term. ``Anchor the
leading term, then process corrections'' is the operational form of
perturbative thinking, and it explains why a model that misses its leading term
fails even when its details are elaborate.

Together the four facets partition the work of closing the round trip:
\emph{where} to expand (prediction), \emph{how} to expand (perturbation),
\emph{what good looks like} (fitness), and \emph{what to hold onto}
(leading-term anchoring).
\section{Two Derivations}
\label{sec:derivations}

From the core definition we derive two structural results. Both concern the
\emph{choice of basis} in which the round trip is expressed.

\subsection{Derivation 1: Basis Choice (why perturbation and hierarchy are equivalent)}
\label{sec:basis}

Perturbation series and hierarchical (level-by-level) structures are both
``leading term plus successive corrections''---the same series structure,
algebraically equivalent. They differ only in the \emph{semantics of the order}:
precision in perturbation, concept in hierarchy. When conceptual crossings
(coupling terms) exist, the equivalence breaks down and a joint basis is
required.

\paragraph{The basis is freely choosable.}
Among perturbation bases, hierarchical bases, tree bases, and Hilbert-space
bases, the optimal basis is the one in which the expansion coefficients are most
sparse (least redundant)---i.e., the most compressed representation. Sparsity of
coefficients is thus the criterion for basis optimality, and it is exactly the
round-trip condition applied to the representation itself.

\paragraph{Hierarchy/tree corresponds strictly to multiresolution bases.}
The rigorous counterpart of hierarchy and tree is multiresolution analysis
(Haar, Daubechies, Faber--Schauder)~\cite{daubechies_1988}: compared with
Fourier bases, multiresolution bases add \emph{locality} (basis functions with
local support) and \emph{nestedness} ($V_j \subset V_{j+1}$, tier $=$ scale).
``Jumping levels'' corresponds to exponential scales $2^{j}$ combined with
tree-like branching. Orthogonality can be preserved (Haar and Daubechies bases
are $L^{2}$-orthogonal), so the round-trip identity is not sacrificed.

\paragraph{Multiplicative structure as tensor products.}
Multiplicative (factorized) structure unifies with the same picture: a single
rooted tree is branching; a product of multiple roots is a tensor product
($1\mathrm{D}\otimes1\mathrm{D}\otimes\cdots$), and the two are unified under
multiresolution analysis. This answers how to describe a factorized structure
that is not a single-root split: as a tensor product of rooted trees, one
factor per root.

\subsection{Derivation 2: Two-Stage Compression (domain selection then refinement; order cannot be reversed)}
\label{sec:twostage}

General (structural) compression is responsible for \emph{selecting the domain};
formula-level compression is responsible for \emph{refining within the domain}
(lossless derivation). The ``losslessness'' of formula-level work is relative to
the domain of applicability, and that domain must be fixed by general thinking.
Reversing the order---refining before selecting the domain---means deriving
precisely in the wrong domain, which drives the result farther off, not closer.

This explains an empirical observation: in some respects formula-level methods
are better while general thinking is not. Formula-level compression excels in
the refinement stage (precise, derivable) but helps little---and can even harm---
in the domain-selection stage, by inducing premature priors.

\paragraph{Domain selection is a three-step process.}
Domain selection is not a static judgment but a process: (i) verify the range in
which the axioms hold; (ii) transfer the axioms that still hold, together with
their conditions; (iii) think about and experiment with new conditions of
applicability. This fills the gap in the reasoning chain between induction and
deduction (induction $\to$ condition verification $\to$ transfer $\to$
experiment $\to$ deduction $\to$ predictive quality), where previously only
deduction was explicit.

\paragraph{Meta-criterion.}
Any mathematical metaphor is itself ``compression plus selection''. Its value is
measured by whether it yields falsifiable predictions: if yes, it is useful; if
no, it is empty compression and should be discarded.
\section{Consistency Test: Seven Thinking-Style Hypotheses}
\label{sec:consistency}

The source corpus contains seven thinking-style hypotheses (paraphrased; exact
wording in Appendix~\ref{app:hypotheses}). We derive each from the framework and
check whether the derivation hits.

\begin{table}[htbp]
\centering
\small
\begin{tabular}{p{4.6cm}p{5.2cm}c}
\toprule
Hypothesis (source corpus) & Derived from framework & Hit \\
\midrule
Escalating style $\to$ stepwise refinement & Perturbation: leading term + successive corrections & \checkmark \\
Negation style $\to$ fitness optimization & Fitness landscape: no binary wrong, continuous fitness & \checkmark \\
Compression judgment $\to$ leading term = most compressed & Basis choice: sparsest expansion = most compressed & \checkmark \\
Cannot compress $=$ do not understand & Round trip fails: leading term not anchored & \checkmark \\
Long time scale $\to$ many perturbation orders & Perturbation convergence: radius requires enough orders & \checkmark \\
Operational theory $\to$ measurable fitness & Fitness: the quantity experiments can measure & \checkmark \\
Sensitive to who is thinking $\to$ AI degradation = fitness drop & Fitness landscape: degraded agent = lower fitness & \checkmark \\
\bottomrule
\end{tabular}
\caption{Seven thinking-style hypotheses derived from the framework. 7/7 hit.
This is an \emph{internal consistency check} on the source corpus---the
framework was constructed from and tested on the same corpus (self-referential)---
not external validation.}
\label{tab:seven}
\end{table}

All seven hypotheses are recovered: 7/7. For comparison, the older
``compression + prediction'' framing recovered only 5/7.

\paragraph{Honesty of the test.}
We state the status of this result explicitly. This is an \emph{internal
consistency check} on the source corpus: the framework was constructed from and
tested on the same corpus, so the test is self-referential. It demonstrates
that the formalization is internally coherent and that its facets are
sufficient to reproduce the observed thinking styles \emph{within} the corpus.
It is \emph{not} external validation: no independent corpus, no held-out
examples, and no cross-subject replication are involved. We claim internal
consistency, not external validity.
\section{Actionable Protocol: Five Injectable Templates}
\label{sec:protocol}

We distill the framework into five templates (T1--T5), each with a trigger
condition and an action. Each template is a short instruction that can be
injected into an agent prompt as a thinking habit. We position them as an
\emph{agent-prompt injection protocol}: the claim is that they are
\emph{injectable and operable} (well-defined triggers, executable actions,
composable). We do \emph{not} claim effectiveness; that is pending independent
experiments.

\begin{enumerate}[leftmargin=2em,itemsep=2pt]
\item \textbf{T1 Cost model.} \emph{Trigger:} a concrete solution has been
found, or ``is $X$ optimal?'' is asked. \emph{Action:} write down the cost model
(what is being costed, under what assumptions) and the unified form; verify the
original solution is a special case; if it cannot be derived, name the missing
assumption.
\item \textbf{T2 Unified coverage.} \emph{Trigger:} facing several seemingly
unrelated methods/systems. \emph{Action:} find one formula or analogy covering
all of them (compression rate / information gain per cost / channel bandwidth),
and write out each term's degeneracy conditions.
\item \textbf{T3 Quantitative landing.} \emph{Trigger:} a qualitative conclusion
is given. \emph{Action:} immediately quantify it into a testable number
(success-rate expression / bandwidth lower bound / expected number of tests); if
it cannot be quantified, mark it ``needs measurement'' and design the cheapest
probe.
\item \textbf{T4 Analogy audit.} \emph{Trigger:} about to borrow a historical
case / analogy / prior. \emph{Action:} check dimension-by-dimension the
condition difference $\Delta$ and the conclusion's sensitivity $s_k$; if
sensitive and mismatched, do not borrow; attach an $\varepsilon\cdot$uniform
floor to prevent confident errors.
\item \textbf{T5 Criterion bottleneck.} \emph{Trigger:} a framework has taken
shape. \emph{Action:} find the ``criterion bottleneck''---which step in the loop
lacks cheap, reliable ground truth; that step is the real bottleneck, not the
architecture. Prefer the remedy ``make the world closed'' (synthetic domains /
executable verification).
\end{enumerate}

\paragraph{Positioning.}
The five templates are the operational layer of the framework: T1--T2 are the
``climb one level'' move (concrete $\to$ cost model $\to$ unified form, original
solution as special case); T3 is predictive power (conclusions land on testable
numbers); T4 is the robustness discipline of leading-term anchoring applied to
borrowed priors; T5 is the fitness-landscape move (locate where ground truth is
missing). Their value proposition is \emph{operability}: each has a crisp
trigger and a mechanical action, so they can be injected into a prompt and
audited for whether they fired. Whether they improve downstream task performance
is an empirical question we leave open.
\section{Related Work}
\label{sec:related}

We position the framework against four adjacent literatures.

\paragraph{Compression as understanding.}
Zipf's law and power laws in cognition treat linguistic and cognitive structure
as evidence of compression-based processing~\cite{newman_powerlaws_2005}. \emph{Our
difference:} we do not fit a power law to data; we define understanding as a
round-trip identity ($\Dmap\circ\Cmap=\mathrm{id}$) and use compression as the
\emph{operational criterion} for when a model is understood, not as a statistical
signature.

\paragraph{Perturbation theory in machine learning.}
Perturbative and influence-based analyses of neural networks characterize how
small changes in parameters or training data propagate to predictions~\cite{kaspschak_nnpt_2020,koh_liang_influence_2017}. \emph{Our difference:}
we use perturbation not to analyze a fixed model but as the \emph{method} by
which a thinker builds understanding term by term, with necessity governed by
round-trip closure (add a correction only when the round trip fails to close).

\paragraph{Multiresolution analysis.}
Multiresolution analysis (Haar, Daubechies, Faber--Schauder) provides nested,
locally supported bases with $V_j\subset V_{j+1}$~\cite{daubechies_1988}.
\emph{Our difference:} we identify multiresolution bases as the rigorous
counterpart of \emph{hierarchical thinking} (tier $=$ scale, jump $=$
exponential scale $2^j$), and use sparsity of coefficients as the criterion for
basis optimality in reasoning, not for signal denoising.

\paragraph{Active learning of specifications / reward hacking as specification discovery.}
Work on reward hacking and specification gaming characterizes misbehavior as a
symptom of underspecified objectives~\cite{skalse_reward_hacking_2022,gallego_spec_self_correction_2025}. \emph{Our
difference:} our T5 template locates the bottleneck as the step lacking cheap,
reliable ground truth, and prescribes ``making the world closed'' (synthetic
domains with known truth) as the primary remedy---a diagnostic-then-remedy
protocol rather than a safety argument.
\section{Limitations}
\label{sec:limitations}

We enumerate the limitations explicitly.

\paragraph{Self-referential consistency test.}
The 7/7 consistency result is an internal consistency check on the source
corpus: the framework was constructed from and tested on the same corpus. It
establishes internal coherence, not external validity. No held-out corpus,
independent subject, or cross-domain replication is involved.

\paragraph{Conceptual framework, not empirical result.}
The framework is a descriptive formalization. It organizes observed thinking
practices into a structure from which they are derivable; it does not measure a
cognitive variable, and it makes no claim about the mechanism of human or AI
thinking.

\paragraph{Template effectiveness untested.}
The five templates (T1--T5) are claimed to be injectable and operable. Whether
injection improves downstream task performance is untested; that is the role of
independent experiments, which we do not perform here.

\paragraph{Single corpus source.}
All structural evidence comes from one source corpus (an anonymized set of
benchmark conversations, from which only structural features are used). The
framework's generality across other corpora, subjects, or domains is an open
question.
\section{Conclusion}
\label{sec:conclusion}

We presented a descriptive formalization of general thinking in which
understanding is a lossless round-trip fixed point of a compression--decompression
pair ($\Dmap\circ\Cmap=\mathrm{id}$, lossless $=$ unique $+$ reproducible). From
this single condition we derived four facets (predictive power, perturbation,
fitness, leading-term anchoring) and two structural results (basis choice via
multiresolution analysis; two-stage domain-selection-then-refinement with
irreversible order). The framework reproduces seven thinking-style hypotheses
with 7/7 internal consistency on the source corpus (an internal consistency
check, not external validation), and we distilled it into five injectable prompt
templates (T1--T5) with explicit triggers and actions.

The contribution is deliberately scoped: a descriptive formalization plus an
actionable, auditable protocol. We do not claim to have discovered the mechanism
of thinking; we claim that a small formalization coherently accounts for the
observed practices \emph{within} the source corpus and yields an injectable
operational layer whose effectiveness awaits independent experiments.

\paragraph{Statement of AI use.}
This manuscript was drafted with AI assistance (large-language-model generation
of the LaTeX draft and formalization) and was reviewed, edited, and approved by
a human author, who takes full responsibility for its content, claims, and
honesty. No AI system is listed as an author.
\appendix
\section{The Seven Hypotheses: Source Corpus Wording and Mapping}
\label{app:hypotheses}

Table~\ref{tab:seven} gives the paraphrase and derivation. Table~\ref{tab:seven_raw}
gives the original (Chinese) wording from the source corpus alongside the
English paraphrase used in the paper, for traceability.

\begin{table}[htbp]
\centering
\small
\begin{tabular}{p{5.4cm}p{5.4cm}}
\toprule
Source corpus wording (original) & English paraphrase (paper) \\
\midrule
升级式$\to$微扰逐级精细化 & Escalating style $\to$ stepwise (perturbative) refinement \\
否定式$\to$适配度优化 & Negation style $\to$ fitness optimization \\
压缩判断$\to$主项=最压缩 & Compression judgment $\to$ leading term $=$ most compressed \\
压不动=没懂$\to$主项没抓稳 & Cannot compress $=$ do not understand $\to$ leading term not anchored \\
时间尺度长$\to$微扰收敛需足够多阶 & Long time scale $\to$ perturbation convergence needs enough orders \\
操作性理论$\to$适配度可测量 & Operational theory $\to$ fitness is measurable \\
对谁在想敏感$\to$AI退化=适配度下降 & Sensitive to who is thinking $\to$ AI degradation $=$ fitness drop \\
\bottomrule
\end{tabular}
\caption{Original wording (source corpus) versus the English paraphrase used in
Section~\ref{sec:consistency}. 7/7 hit, internal consistency only.}
\label{tab:seven_raw}
\end{table}

\section{Number Traceability}
\label{app:trace}

Every quantitative statement in this paper maps to the source corpus as follows
(the accompanying \texttt{number-trace.md} gives file paths and line numbers):

\begin{itemize}
\item \textbf{7/7 consistency} --- the seven hypotheses and the 7/7 (vs.\ 5/7
for the older ``compression + prediction'' framing) are recorded in the
``inspection'' section of the core-framework source file.
\item \textbf{Five templates (T1--T5)} --- the five injectable templates with
triggers and actions are recorded in the ``injectable templates'' section of the
template source file.
\item \textbf{Four facets} --- the four facets (predictive power, perturbation,
fitness, leading-term anchoring) are recorded in the ``facets'' section of the
core-framework source file.
\item \textbf{Two derivations} --- basis choice (multiresolution bases) and
two-stage compression (domain selection $\to$ refinement) are recorded in the
two ``derivation'' sections of the core-framework source file.
\item \textbf{Three-step domain selection} --- the three-step process
(axiom-range verification, conditional transfer, new-condition experimentation)
is recorded in the two-stage-compression derivation section.
\end{itemize}

The source corpus is private; only structural features and the paraphrased
hypotheses above are used. No verbatim benchmark-conversation content appears in
this paper.

\begin{thebibliography}{9}
\bibitem{newman_powerlaws_2005}
M.~E.~J. Newman. Power laws, Pareto distributions and Zipf's law.
\emph{Contemporary Physics}, 46(5):323--351, 2005. doi:10.1080/00107510500052444.
\bibitem{kaspschak_nnpt_2020}
B.~Kaspschak and U.-G. Meißner. A neural network perturbation theory based on
the Born series. arXiv:2009.03192, 2020.
\bibitem{koh_liang_influence_2017}
P.~W. Koh and P.~Liang. Understanding black-box predictions via influence
functions. arXiv:1703.04730, 2017.
\bibitem{daubechies_1988}
I.~Daubechies. Orthonormal bases of compactly supported wavelets.
\emph{Communications on Pure and Applied Mathematics}, 41(7):909--999, 1988.
doi:10.1002/cpa.3160410705.
\bibitem{skalse_reward_hacking_2022}
J.~Skalse, N.~H.~R. Howe, D.~Krasheninnikov, and D.~Krueger. Defining and
characterizing reward hacking. arXiv:2209.13085, 2022.
\bibitem{gallego_spec_self_correction_2025}
V.~Gallego et al. Specification self-correction: mitigating in-context reward
hacking through test-time refinement. arXiv:2507.18742, 2025.
\end{thebibliography}

\end{document}
